% =====================================================================
%  CAPÍTULO 3: DESARROLLO DE LA PROPUESTA
% =====================================================================
\chapter{Desarrollo de la Propuesta}

\section{Metodología de Desarrollo}
El modelo se desarrollará siguiendo CRISP-DM \citep{chapman2000}. Como se trabaja con relatos sensibles, a sus etapas se les añaden la autorización ética, la anotación y la documentación. Para que el trabajo responda directamente a los objetivos, las seis etapas de CRISP-DM se agrupan en cuatro etapas de desarrollo, una por cada objetivo específico (Tabla~\ref{tab:crisp}).

\begin{table}[H]
\begin{threeparttable}
\caption{Relación entre objetivos específicos, etapas de desarrollo y CRISP-DM}
\label{tab:crisp}
\singlespacing\small
\renewcommand{\arraystretch}{1.25}
\begin{tabular}{@{}L{1.6cm} L{4.2cm} L{3.9cm} L{5.7cm}@{}}
\toprule
\textbf{Objetivo} & \textbf{Etapa de desarrollo} & \textbf{Etapas de CRISP-DM} & \textbf{Producto} \\
\midrule
OE1 & Etapa 1: Construcción del corpus & Comprensión del problema, comprensión y preparación de los datos & Corpus anonimizado y anotado, con su kappa y su hoja de datos \\
OE2 & Etapa 2: Entrenamiento y selección del modelo & Preparación de los datos y modelado & Modelo seleccionado y registro de experimentos \\
OE3 & Etapa 3: Validación del modelo & Evaluación & Métricas en la prueba, contraste de hipótesis y análisis de errores \\
OE4 & Etapa 4: Implementación del prototipo & Despliegue & Prototipo verificado y ficha técnica del modelo \\
\bottomrule
\end{tabular}
\begin{tablenotes}[flushleft]
\small
\item \textit{Nota.} Adaptado de \citet{chapman2000}.
\end{tablenotes}
\end{threeparttable}
\end{table}

\section{Arquitectura de la Solución}
La solución tiene dos flujos, que se muestran en la Figura~\ref{fig:arquitectura}. En el flujo de entrenamiento, el corpus anotado se convierte en un modelo validado. En el flujo de inferencia, el prototipo usa ese modelo para analizar un relato nuevo.

\begin{figure}[H]
\caption{Arquitectura de la solución}
\label{fig:arquitectura}
\centering
\begin{tikzpicture}[
	node distance=0.5cm and 0.45cm,
	caja/.style={draw, rounded corners=3pt, fill=diagAzul!10, align=center,
	text width=2.55cm, minimum height=1.25cm, font=\footnotesize},
	caja2/.style={caja, fill=diagVerde!15},
	flecha/.style={-{Stealth[length=2.2mm]}, thick},
	etiqueta/.style={font=\footnotesize\bfseries, anchor=west}
	]
	\node[caja] (c1) {Corpus anonimizado y anotado};
	\node[caja, right=of c1] (c2) {Preprocesamiento y partición por participante};
	\node[caja, right=of c2] (c3) {Representación\\TF-IDF / BETO};
	\node[caja, right=of c3] (c4) {Entrenamiento, selección y validación};
	\node[caja, below=0.9cm of c4, fill=diagRojo!12] (m) {\textbf{Modelo validado}};
	\node[caja2, below=2.8cm of c1] (i1) {Relato nuevo};
	\node[caja2, right=of i1] (i2) {Revisión de datos personales};
	\node[caja2, right=of i2] (i3) {Predicción y probabilidad};
	\node[caja2, right=of i3] (i4) {Explicación SHAP y revisión humana};
	\draw[flecha] (c1) -- (c2);
	\draw[flecha] (c2) -- (c3);
	\draw[flecha] (c3) -- (c4);
	\draw[flecha] (c4) -- (m);
	\draw[flecha] (i1) -- (i2);
	\draw[flecha] (i2) -- (i3);
	\draw[flecha] (i3) -- (i4);
	\draw[flecha, dashed] (m) -| (i3);
	\node[etiqueta] at ([yshift=0.35cm]c1.north west) {Flujo de entrenamiento};
	\node[etiqueta] at ([yshift=0.35cm]i1.north west) {Flujo de inferencia (prototipo)};
\end{tikzpicture}

\raggedright\textit{Nota.} Elaboración propia.
\end{figure}

\section{Etapas del Desarrollo}
La Figura~\ref{fig:fases} resume las cuatro etapas. Cada una corresponde a un objetivo específico y entrega un producto que la siguiente necesita.

\begin{figure}[H]
\caption{Etapas del desarrollo del modelo}
\label{fig:fases}
\centering
\begin{tikzpicture}[
	node distance=0.6cm,
	fase/.style={draw, rounded corners=3pt, fill=diagAzul!10, align=center,
	text width=3.2cm, minimum height=1.9cm, font=\small},
	flecha/.style={-{Stealth[length=2.5mm]}, thick}
	]
	\node[fase] (e1) {\textbf{Etapa 1 (OE1)}\\Construcción del corpus\\[2pt]\footnotesize Autorización, piloto, recolección, anonimización y anotación};
	\node[fase, right=of e1] (e2) {\textbf{Etapa 2 (OE2)}\\Entrenamiento y selección\\[2pt]\footnotesize Partición, 13 configuraciones y validación cruzada};
	\node[fase, right=of e2] (e3) {\textbf{Etapa 3 (OE3)}\\Validación\\[2pt]\footnotesize Prueba única, hipótesis y análisis de errores};
	\node[fase, right=of e3, fill=diagVerde!15] (e4) {\textbf{Etapa 4 (OE4)}\\Prototipo\\[2pt]\footnotesize Servidor, interfaz, pruebas y documentación};
	\draw[flecha] (e1) -- (e2);
	\draw[flecha] (e2) -- (e3);
	\draw[flecha] (e3) -- (e4);
\end{tikzpicture}

\raggedright\textit{Nota.} Elaboración propia.
\end{figure}

\subsection{Etapa 1: Construcción del corpus (OE1)}
Primero se define para qué se usará el modelo y quién custodiará los datos, y se obtiene la aprobación del comité de ética y de la universidad. Los especialistas revisan la guía de anotación antes de empezar.

Luego se aplica el formulario a unos 40 estudiantes para comprobar que las preguntas se entienden y que la anonimización y la anotación funcionan. Con los ajustes necesarios se abre la convocatoria principal, siempre voluntaria y sin acceder a conversaciones privadas ni a expedientes.

A cada relato recibido se le quitan nombres, teléfonos, correos, códigos y cualquier detalle que permita reconocer a alguien, y se reemplazan por marcadores como \texttt{[PERSONA]} o \texttt{[CURSO]}. Después, dos anotadores etiquetan cada relato por separado y se calcula su acuerdo con el kappa de Cohen (Ecuación~\ref{eq:kappa}); un especialista resuelve los casos en que no coinciden. La etapa termina con el corpus anotado y su hoja de datos \citep{gebru2021}.

\subsection{Etapa 2: Entrenamiento y selección del modelo (OE2)}
Se reserva alrededor del 20\,\% de los participantes para la prueba final y el 80\,\% restante se usa para el desarrollo, manteniendo la proporción de relatos con y sin acoso. Antes de entrenar cualquier modelo se define la lista de palabras de la línea base léxica (Ecuación~\ref{eq:lexica}).

Con los datos de desarrollo se entrenan las 13 configuraciones del diseño. Los relatos se representan con TF-IDF (Ecuación~\ref{eq:tfidf}) para la regresión logística, la SVM lineal y Naïve Bayes (Ecuaciones~\ref{eq:logistica} a \ref{eq:nb}), y con la representación contextual de BETO (Ecuaciones~\ref{eq:atencion} y \ref{eq:cls}). Cada modelo se prueba sin tratamiento del desbalance, con ponderación de clases (Ecuación~\ref{eq:pesos}) y, en el caso de TF-IDF, con SMOTE (Ecuación~\ref{eq:smote}). Los hiperparámetros se ajustan con validación cruzada agrupada de cinco pliegues. Se elige el modelo con mayor macro F1; si dos empatan, se prefiere el de mayor F1 en la clase positiva y, después, el más sencillo. El Algoritmo~\ref{alg:entrenamiento} resume el procedimiento.

\begin{algorithm}[H]
\caption{Entrenamiento, selección y validación del modelo}
\label{alg:entrenamiento}
\singlespacing
\begin{algorithmic}[1]
	\Require Corpus anotado $D$ con etiqueta e identificador de participante; configuraciones $\mathcal{C}$
	\Ensure Modelo seleccionado $m^{*}$ y sus métricas en la prueba
	\State Agrupar relatos duplicados o casi idénticos
	\State Dividir $D$ por participante en $D_{dev}$ (80\,\%) y $D_{test}$ (20\,\%) \Comment{$D_{test}$ no se toca}
	\State Crear 5 pliegues por participante dentro de $D_{dev}$
	\ForAll{$c \in \mathcal{C}$ y las dos líneas base}
		\ForAll{pliegue $k$}
			\State Ajustar la representación y el remuestreo \textbf{solo} con el entrenamiento del pliegue $k$
			\State Entrenar $c$ y registrar macro F1 y F1 positivo en la validación del pliegue $k$
		\EndFor
		\State Calcular el promedio y la dispersión de las métricas de $c$
	\EndFor
	\State $m^{*} \gets$ configuración con mayor macro F1
	\State Reentrenar $m^{*}$ con todo $D_{dev}$
	\State Evaluar una sola vez en $D_{test}$: $P$, $R$, $F_1$, macro $F_1$ y matriz de confusión
	\State \Return $m^{*}$ y sus métricas con intervalos del 95\,\%
\end{algorithmic}
\end{algorithm}

\subsection{Etapa 3: Validación del modelo (OE3)}
El modelo elegido y las líneas base se prueban una sola vez con los relatos reservados. Se calculan la precisión, la exhaustividad, el F1 y el macro F1 (Ecuaciones~\ref{eq:metricas} y \ref{eq:macrof1}), con sus intervalos de confianza por \textit{bootstrap}, y se contrastan las hipótesis con la prueba de McNemar (Ecuación~\ref{eq:mcnemar}) según los criterios de la Tabla~\ref{tab:contraste}. Después se revisan uno por uno los relatos en los que el modelo se equivocó, para entender por qué: negaciones, citas, insinuaciones, falta de contexto, expresiones locales o ironía.

\subsection{Etapa 4: Implementación del prototipo (OE4)}
Se implementa el prototipo descrito en la siguiente sección: el servidor con FastAPI que carga el modelo validado y calcula las explicaciones con SHAP (Ecuación~\ref{eq:shap}), y la interfaz en React. Luego se verifica con relatos de ejemplo que el prototipo dé las mismas predicciones que el modelo y que cumpla cada requisito funcional. El modelo se documenta con una ficha técnica \citep{mitchell2019} y, por último, se vuelve a ejecutar todo el proceso desde cero para comprobar que se obtienen las mismas métricas.

\section{Prototipo de Apoyo a la Revisión}
El prototipo será una aplicación web que funciona en un equipo local. Tendrá un servicio en Python con FastAPI, que carga el modelo validado, y una interfaz en React para quien revisa los relatos. El usuario pega un relato, el sistema le avisa si detecta posibles datos personales y luego muestra si el relato describe acoso, con qué probabilidad y qué palabras influyeron más en esa decisión según SHAP \citep{lundberg2017}. Los relatos analizados no se guardan ni se envían a ningún servicio externo. Las Tablas~\ref{tab:rf} y \ref{tab:rnf} listan sus requisitos.

\begin{table}[H]
\begin{threeparttable}
\caption{Requisitos funcionales del prototipo}
\label{tab:rf}
\singlespacing\small
\renewcommand{\arraystretch}{1.25}
\begin{tabular}{@{}L{1.5cm} L{13.9cm}@{}}
\toprule
\textbf{Código} & \textbf{Requisito} \\
\midrule
RF01 & Recibir un relato en español a través de un campo de texto. \\
RF02 & Avisar si el relato contiene posibles datos personales, como nombres, correos, teléfonos o códigos. \\
RF03 & Clasificar el relato como \textit{con acoso sexual} o \textit{sin acoso sexual} con el modelo validado. \\
RF04 & Mostrar la probabilidad de que el relato describa acoso sexual. \\
RF05 & Mostrar las palabras que más aumentan y disminuyen esa probabilidad. \\
RF06 & Mostrar siempre un aviso de que el resultado es un apoyo a la revisión y no una decisión. \\
\bottomrule
\end{tabular}
\begin{tablenotes}[flushleft]
\small
\item \textit{Nota.} Elaboración propia.
\end{tablenotes}
\end{threeparttable}
\end{table}

\begin{table}[H]
\begin{threeparttable}
\caption{Requisitos no funcionales del prototipo}
\label{tab:rnf}
\singlespacing\small
\renewcommand{\arraystretch}{1.25}
\begin{tabular}{@{}L{1.5cm} L{13.9cm}@{}}
\toprule
\textbf{Código} & \textbf{Requisito} \\
\midrule
RNF01 & Privacidad: procesar los relatos en el equipo local, sin guardarlos ni enviarlos fuera. \\
RNF02 & Fidelidad: dar las mismas predicciones que el modelo validado. \\
RNF03 & Rendimiento: responder lo bastante rápido para usarse en una revisión interactiva en una computadora común. \\
RNF04 & Mantenibilidad: código versionado, con las dependencias y la versión del modelo registradas. \\
\bottomrule
\end{tabular}
\begin{tablenotes}[flushleft]
\small
\item \textit{Nota.} Elaboración propia.
\end{tablenotes}
\end{threeparttable}
\end{table}

\section{Resultados Esperados}
Al terminar la investigación se espera contar con:
\begin{itemize}
	\item Una guía de anotación revisada por especialistas y un informe del acuerdo entre anotadores.
	\item Un corpus anonimizado de relatos, de acceso restringido, con su hoja de datos.
	\item Un modelo de aprendizaje automático que detecte acoso sexual mejor que las líneas base, con su ficha técnica y el análisis de sus errores.
	\item El código del proceso completo, versionado y acompañado de datos de ejemplo.
	\item Un prototipo funcional que apoye a las oficinas de orientación y atención de la universidad en la revisión de relatos, siempre bajo el criterio de su personal y dentro de los procedimientos que establece la normativa \citep{ds0142019}.
\end{itemize}
